% ===================================================================
%  TAULA N° 8 — La Meissa. — La Caça, la Vendémia, la Pesca.
%  Traduction 2026 d'apres texto/io, controlee sur texto/fr et
%  texto/ca. Consigne : voir texto/oc/10-taula-01.tex.
%
%  L'appel de note est dans le TITRE de la scene. LE RENVOI (13) DE
%  LA FAUX MANQUE A L'IDO ; la colonne le rend : ecart declare.
% ===================================================================

%%K t08-noto-1 noto
\VUnotes{143.2mm}{%
\textit{(*) Qu'aqueste mot es pas precís: s'agís d'una meissa de lin, d'una
vendémia, d'una culhida de poms? Benlèu que valdriá mai emplegar «}
\VUgras{mesono} \textit{», prepausat dins} Mondo \textit{XIV, p. 242. Mas fins
ara aquela raiç novèla es pas estada adoptada per l'Acadèmia e espèra la
discussion segon las règlas costumièras del movement idista.}%
}

%%K t08-apar-1 apar
\VUsaut{5.0mm}
\VUtitre{15.0pt}{60}{TAULA N° 8}
\VUsaut{3.0mm}
\VUcentre{13.2pt}{90}{\textit{Primièra scèna.}}
\VUsaut{3.0mm}
\VUpk{10.2pt}{40}{La Meissa (*).}
\VUsaut{4.0mm}
\VUpk{10.2pt}{40}{Los Aspèctes de la Campanha.}
\VUsaut{3.0mm}

%%K t08-c1-01-1 p
1. --- Amb l'\VUgras{estiu}, la campanha a cambiat d'aspècte un còp de mai. La
grana confisada al rèc a germinat; una planteta verda es sortida de la tèrra e a
donat son espiga. Lo gran s'es format, puèi a madurat, e demòra pas mai que de
meissonar.

%%K t08-c1-02-1 p
2. --- Las \VUgras{flors dels camps} son dobèrtas. \VUgras{Blavets}
\textsuperscript{(1)}, \VUgras{rosèlas} \textsuperscript{(2)} e \VUgras{didalièras}
\textsuperscript{(3)} pausan dins lo ton uniforme dels blats l'alègre contrast de
lors frescas \VUgras{coròlas}.

%%K t08-c1-03-1 p
3. --- Los arbres fruchièrs se son cobèrts de fuèlhas; lo fruch a madurat. Lo
pichon \VUgras{cerièr} \textsuperscript{(4)} es cargat de bèlas \VUgras{cerièras}
\textsuperscript{(5)}, que los enfants culhisson e manjan amb delici.

%%K t08-c1-04-1 p
4. --- Lo tropèl pòt ja passar tot lo jorn a l'aire liure.

%%K t08-c1-04-2 p
Los \VUgras{anhèls} \textsuperscript{(6)} an creissut e son venguts de bèlas
\VUgras{fedas} \textsuperscript{(7)} e de bèles \VUgras{marrans}
\textsuperscript{(8)} amb un velon espés. Pasturan tranquillament l'\VUgras{èrba}
del \VUgras{pastenc} \textsuperscript{(9)}, esposcat de \VUgras{margaridas}.

%%K t08-c1-04-3 p
Lèu començaràn de tondre aqueles animals per n'aver la \VUgras{lana}, amb la
quala se fabrica lo drap de nòstra rauba.

%%K t08-tit-2 sub
\VUsaut{5.0mm}
\VUpk{10.2pt}{40}{Lo Pastre.}
\VUsaut{3.4mm}

%%K t08-c1-05-1 p
5. --- Lo \VUgras{pastre} \textsuperscript{(10)} sembla pas ne far gaire cas.
Se'n sosca pas que de l'auratge que passa a l'orizont, e lo \VUgras{pastorèl}
\textsuperscript{(13)}, que s'apròcha la \VUgras{gorda} \textsuperscript{(14)} de
las pòtas, sembla encara mai despreocupat.

%%K t08-c1-06-1 p
6. --- La calor es aclapanta; que lo \VUgras{gos} \textsuperscript{(15)} ven
tanben se refrescar; lapa l'aiga fresca del \VUgras{riu} \textsuperscript{(16)} e
espauruga una paura \VUgras{granolha} \textsuperscript{(17)} qu'èra amagada dins
l'èrba de la riba. Aquesta sauta dins l'aiga clara ont florisson los
\VUgras{nenufars} \textsuperscript{(18)}.

%%K t08-c1-07-1 p
7. --- Una \VUgras{coalonga} \textsuperscript{(19)} se banha contra la pichona
\VUgras{font} \textsuperscript{(20)} que raja entre doas ròcas e s'amaga jos la
\VUgras{bartassa espinosa} \textsuperscript{(21)} e los \VUgras{bartasses
d'albespin} \textsuperscript{(22)}.

%%K t08-tit-3 sub
\VUsaut{5.0mm}
\VUpk{10.2pt}{40}{La Meissa.}
\VUsaut{3.4mm}

%%K t08-c1-08-1 p
8. --- Pels enfants de la campanha, la sèga del blat es una epòca fòrça
divertenta. Fan d'alègras \VUgras{cambaviradas} \textsuperscript{(24)} suls
\VUgras{molons de palha} \textsuperscript{(23)}, ajudan los \VUgras{segaires}
\textsuperscript{(26)} e, mentre qu'aqueles sègan lo \VUgras{camp de blat}
\textsuperscript{(27)} amb lors \VUgras{dalhas} \textsuperscript{(25)}, plan
agusadas a la \VUgras{pèira d'agusar} \textsuperscript{(30)}, eles recampan lo
blat e lo ligan en \VUgras{garbas} \textsuperscript{(31)} o en \VUgras{fais}
\textsuperscript{(32)}.

%%K t08-c1-09-1 p
9. --- De còps se permet als grands de talhar amb un \VUgras{volam}
\textsuperscript{(12)} las \VUgras{tijas dauradas} \textsuperscript{(28)} que
pòrtan las pesucas \VUgras{espigas} \textsuperscript{(29)} plenas de gran; e los
pichons, mentre que gardan lo \VUgras{bestial} \textsuperscript{(33)} que pastura
dins lo \VUgras{prat} \textsuperscript{(34)} a l'ombra del grand \VUgras{telh}
\textsuperscript{(35)}, caçan las bèlas \VUgras{parpalhòlas} amb lor
\VUgras{filat} \textsuperscript{(36)}.

%%K t08-c1-09-2 p
D'unes pugen quitament sul \VUgras{carri} \textsuperscript{(37)} per pausar las
garbas que lor passan amb una \VUgras{forca} \textsuperscript{(38)}.

%%K t08-c1-10-1 p
10. --- Per tota la \VUgras{plana} \textsuperscript{(39)}, ont s'enauran las
\VUgras{alausas} \textsuperscript{(41)}, i règna l'animacion. Totes son ocupats
amb la sèga. Mas, mentre que los paures contúnhan d'emplegar los espleches ancians
--- \VUgras{dalhas, volams} e \VUgras{flagèls} \textsuperscript{(40)} ---, los
proprietaris rics fan sègar lor blat o lor \VUgras{segal} \textsuperscript{(43)}
amb una \VUgras{segaira mecanica} \textsuperscript{(42)}. Puèi, sens besonh de
portar las garbas al granièr, se fan aquí meteis de grands \VUgras{molons de
garbas} \textsuperscript{(44)}.

%%K t08-c1-11-1 p
11. --- Alara se pòrta una \VUgras{batusa} \textsuperscript{(47)} qu'una
\VUgras{locomobila} \textsuperscript{(45)} mòu per mejan d'una \VUgras{correja de
transmission} \textsuperscript{(46)}, e atal lo gran demòra separat de la
\VUgras{palha} sens sortir del camp ont foguèt semenat.

%%K t08-tit-4 sub
\VUsaut{5.0mm}
\VUpk{10.2pt}{40}{L'Auratge.}
\VUsaut{3.4mm}

%%K t08-c1-12-1 p
12. --- Sovent esclatan de violents \VUgras{auratges} \textsuperscript{(53)} a
l'epòca de la sèga. D'en primièr s'ausisson de luènh los ronsaments del
\VUgras{tròn}; se lèva un fòrt \VUgras{vent de ponent} \textsuperscript{(54)} que
plega amb de gemècs los arbres mai gròsses del \VUgras{bòsc}
\textsuperscript{(49)}. De negres \VUgras{nivolasses} \textsuperscript{(56)}
assautan lo cèl; los traucan los zigazagas encegants de l'\VUgras{esluc}
\textsuperscript{(55)}. La \VUgras{pluèja} e de còps la \VUgras{granissa}
començan de tombar, en aclapant la meissa prèsta a èsser segada.

%%K t08-c1-13-1 p
13. --- Sovent lo fòldre aluca quauque bastiment. Atal, l'an passat lo teulat del
\VUgras{molin de vent} \textsuperscript{(50)} demorèt reduch en cendres e doas de
sas \VUgras{alas} \textsuperscript{(51)} demorèron destrusidas.

%%K t08-c1-14-1 p
14. --- Al cap de qualques oras torna la calma. Un brilhant \VUgras{arc de Sant
Martin} \textsuperscript{(52)} apareis a l'orizont e la natura repren sa vida,
interrompuda un instant. Los païsans, que s'èran refugiats dins las
\VUgras{cabanas} \textsuperscript{(48)}, tòrnan al trabalh e s'esfòrçan valentament
d'adobar los damatges que venon de sofrir.

%%K t08-tit-5 sub
\VUsaut{6.0mm}
\VUcentre{13.2pt}{90}{\textit{Segonda scèna.}}
\VUsaut{3.6mm}

%%K t08-tit-6 sub
\VUpk{10.2pt}{40}{La Caça. --- La Vendémia.}
\VUsaut{1.4mm}
\VUpk{10.2pt}{40}{La Pesca.}
\VUsaut{4.0mm}

%%K t08-c2-01-1 p
1. --- Cap a la fin d'\VUgras{agost} o a mièg \VUgras{setembre}, segon las
regions, se dobrís la caça. Tot bèl just los primièrs rais del solelh an fach
sortir las \VUgras{lagramusas} \textsuperscript{(2)} de son trauc, que ja lo
\VUgras{caçaire} \textsuperscript{(5)} se met en camin amb sos \VUgras{gosses}
\textsuperscript{(4)}.

%%K t08-c2-02-1 p
2. --- S'es pausat son solid \VUgras{vestit de caça} \textsuperscript{(9)} de drap
gròs e de \VUgras{gèstras} de cuèr amb las qualas poirà caminar dins l'èrba mòlha
ont s'amagan los \VUgras{grapauds} \textsuperscript{(1)}; a pres son
\VUgras{escopeta} \textsuperscript{(6)} e son \VUgras{carnièr}
\textsuperscript{(10)}, e ja se n'es anat.

%%K t08-c2-03-1 p
3. --- L'ardor de la caça lo mena al mitan d'una vinha ont la vendémia es en plen
trin. Los filhs del vinhairon, que gardan de costuma las cabras al camin, uèi las
daissan pasturar a lor grat e, aprèp aver ligat a una estaca un vièlh
\VUgras{boc} \textsuperscript{(3)} que mancariá pas de profitar de sa libertat per
s'escapar, se donan tanben lor part de plaser. Un pren un rasim dins una
\VUgras{semala} \textsuperscript{(12)} e se divertís de far córrer lo \VUgras{suc}
\textsuperscript{(14)} dins un \VUgras{veire} \textsuperscript{(13)} per far de
most (vin sens fermentar). L'autre se corona d'una \VUgras{garlanda de fuèlhas}
\textsuperscript{(15)} que ven de trenar. Contra eles, de \VUgras{passeròtas}
\textsuperscript{(16)} descaradas arriban fins al truèlh per pigotar los rasims.

%%K t08-c2-04-1 p
4. --- Mentre que los enfants se divertisson, los òmes trabalhan. Los
\VUgras{vendemiaires} \textsuperscript{(18)} pòrtan lo rasim a la cava dins de
\VUgras{desquets} \textsuperscript{(17)}; un \VUgras{botièr}
\textsuperscript{(19)} adoba las \VUgras{barricas} \textsuperscript{(20)},
verifica se las \VUgras{doguas} \textsuperscript{(22)} e los \VUgras{fonses}
\textsuperscript{(21)} son en bon estat e cambia los \VUgras{cèrcles}
\textsuperscript{(23)} copats. Es tanben el que faguèt las grandas
\VUgras{tinas} \textsuperscript{(25)} ont se voida lo \VUgras{rasim}
\textsuperscript{(26)} per que fermente.

%%K t08-c2-05-1 p
5. --- Quand la fermentacion es acabada, se tira lo vin e se pausa la rèsta al
\VUgras{truèlh} \textsuperscript{(28)} e, en sarrar a plan plan lo grand
\VUgras{cargòl} \textsuperscript{(29)}, se sarra lo suc, que cor dins una
\VUgras{tineta} \textsuperscript{(32)}, e s'obten encara de \VUgras{vin}
\textsuperscript{(31)}.

%%K t08-tit-8 sub
\VUsaut{6.0mm}
\VUpk{10.2pt}{40}{Cervesa e Sidra.}
\VUsaut{3.4mm}

%%K t08-c2-06-1 p
6. --- Per la paret de la \VUgras{cava} \textsuperscript{(27)} s'enfila una branca
de \VUgras{lupol} \textsuperscript{(30)}. Aicí es pas qu'una planta d'ornament, mas
dins d'unes païses, ont la vinha es fòrça rara, lo lupol se cultiva per fabricar
la \VUgras{cervesa}.

%%K t08-c2-07-1 p
7. --- Lo pichon \VUgras{pomièr} \textsuperscript{(33)} es cargat de poms que,
quand seràn madurs, donaràn una \VUgras{sidra} excellenta. Dins sa \VUgras{gàbia}
\textsuperscript{(34)}, lo \VUgras{canari} \textsuperscript{(35)} e lo
\VUgras{cardin} \textsuperscript{(36)} agachan amb enveja las passeròtas que
folastrejan en libertat.

%%K t08-c2-08-1 p
8. --- Tant luènh que la vista pòrta, sul pendís dels truques, s'alinhan los
\VUgras{escassièrs} \textsuperscript{(40)} que sostenon las magnificas
\VUgras{socas} \textsuperscript{(39)}, cargadas de pesucs \VUgras{rasims}.

%%K t08-c2-08-2 p
Pendent tot l'an, lo \VUgras{vinhairon} \textsuperscript{(42)} a trabalhat per
suènhar sa \VUgras{vinha} \textsuperscript{(38)}, e ara es recompensat de son
esfòrç per una bèla meissa. Las \VUgras{vendemiairas} \textsuperscript{(41)}
emplenan las tinas pausadas sul carri que tiran las \VUgras{muòlas}
\textsuperscript{(43)}, e tot lo mond es alègre.

%%K t08-tit-9 sub
\VUsaut{5.0mm}
\VUpk{10.2pt}{40}{La Pesca.}
\VUsaut{3.4mm}

%%K t08-c2-09-1 p
9. --- Es lo meteis pels \VUgras{pescaires a la cana} \textsuperscript{(44)}, que
dempuèi lo matin son venguts s'installar a la riba de l'aiga amb lors
\VUgras{canas de pescar} \textsuperscript{(45)}.

%%K t08-c2-09-2 p
Lo \VUgras{peis} \textsuperscript{(47)} mòrd a l'ime tanlèu que getan lor
\VUgras{fial} \textsuperscript{(46)} dins l'aiga, e an tot bèl just lo temps de
renovelar l'\VUgras{esca} que ja una novèla victima coeja al cap del fial.

%%K t08-c2-09-3 p
Çaquelà, lo \VUgras{pescaire al filat} \textsuperscript{(48)}, que de sa
\VUgras{barca} \textsuperscript{(49)} geta l'esparvièr al mitan del \VUgras{riu}
\textsuperscript{(50)}, ne pren fòrça mai.

%%K t08-c2-10-1 p
10. --- L'auton es la sason dels bons passejaments: a pè, a \VUgras{caval}
\textsuperscript{(52)}, en \VUgras{bicicleta} \textsuperscript{(53)}, en
\VUgras{automobila} \textsuperscript{(54)}. Los \VUgras{camins}
\textsuperscript{(51)} son nets e se profitan los darrièrs bons jorns, que ja las
\VUgras{dindoletas} \textsuperscript{(56)} s'acampan suls teulats per volar a tota
ala cap a de païses mai cauds. Las fuèlhas del vièlh \VUgras{noguièr}
\textsuperscript{(57)} jaunisson, las \VUgras{noses} tomban e los auts
\VUgras{arbres} acampats en \VUgras{boscatge} \textsuperscript{(58)} sus
l'\VUgras{autura} \textsuperscript{(59)} demoraràn lèu sens fuèlhas.

%%K t08-c2-11-1 p
11. --- Lo \VUgras{solelh} \textsuperscript{(60)} se lèva mai tard, los jorns
demenisson, comença de se sentir un freg pro viu al \VUgras{matin}, e ja de
bandas de \VUgras{becassas} \textsuperscript{(61)} daissan nòstres climas pas
ospitalièrs.
\VUsaut{6.0mm}
\VUfilet{22mm}
